\chapter{Demostraciones}
\label{anx:demostraciones}

Este anexo reúne las demostraciones y los desarrollos de los resultados enunciados en el capítulo~\ref{cap:strata}. Cada sección recoge la prueba o la derivación de un resultado, indicado al comienzo por su enunciado y su label en el cuerpo de la memoria.

\section{Recursión forward del HMM}
\label{anx:forward}

Desarrollo de la Proposición~\ref{prop:forward} (Recursión forward), Sección~\ref{sec:hmm-forward}; seguimos el tratamiento estándar del HMM \cite{rabiner1989, bishop2006}.

\begin{proof}
La inicialización es inmediata: $\alpha_1(k) = \Prob(x_1, z_1 = k) = \Prob(z_1 = k)\,\Prob(x_1 \mid z_1 = k) = \pi_k\, b_k(x_1)$, usando $(\ast)$ en el segundo factor.

Para el paso recursivo partimos de la definición e introducimos el estado anterior $z_t$ marginalizando sobre sus $K$ valores posibles:
\[
\alpha_{t+1}(j)
= \Prob(x_{1:t+1},\, z_{t+1} = j)
= \sum_{k=1}^{K} \Prob(x_{1:t+1},\, z_t = k,\, z_{t+1} = j).
\]
Descomponemos cada sumando separando la última observación $x_{t+1}$ del resto mediante la regla del producto,
\[
\Prob(x_{1:t+1}, z_t = k, z_{t+1} = j)
= \Prob(x_{t+1} \mid x_{1:t}, z_t = k, z_{t+1} = j)\,
  \Prob(x_{1:t}, z_t = k, z_{t+1} = j).
\]
Por la independencia condicional de las emisiones $(\ast)$, la observación $x_{t+1}$ depende solo de su propio estado $z_{t+1} = j$, de modo que el primer factor es
\[
\Prob(x_{t+1} \mid x_{1:t}, z_t = k, z_{t+1} = j) = \Prob(x_{t+1} \mid z_{t+1} = j) = b_j(x_{t+1}).
\]
El segundo factor se descompone a su vez aislando la transición $z_t = k \to z_{t+1} = j$:
\[
\Prob(x_{1:t}, z_t = k, z_{t+1} = j)
= \Prob(z_{t+1} = j \mid x_{1:t}, z_t = k)\, \Prob(x_{1:t}, z_t = k).
\]
La Markovianidad de la cadena, junto con la independencia condicional de las emisiones $(\ast)$, hace que, fijado $z_t$, el estado $z_{t+1}$ sea independiente de $x_{1:t}$ y de los estados anteriores (es la d-separación en el grafo del modelo); por tanto $\Prob(z_{t+1} = j \mid x_{1:t}, z_t = k) = A_{kj}$. El factor restante es, por definición, $\alpha_t(k)$. Reuniendo los tres trozos,
\[
\Prob(x_{1:t+1}, z_t = k, z_{t+1} = j) = b_j(x_{t+1})\, A_{kj}\, \alpha_t(k),
\]
y al sumar en $k$ el factor $b_j(x_{t+1})$ no depende del índice de la suma y sale fuera:
\[
\alpha_{t+1}(j) = b_j(x_{t+1}) \sum_{k=1}^{K} \alpha_t(k)\, A_{kj},
\]
que es lo que se quería probar.
\end{proof}

\section{Filtrado y causalidad del HMM}
\label{anx:filtrado}

Desarrollo de la Proposición~\ref{prop:filtrado} (Filtrado a partir del forward) y del Corolario~\ref{cor:causal} (Causalidad del filtrado), Sección~\ref{sec:hmm-filtrado}, según el tratamiento estándar del HMM \cite{rabiner1989, bishop2006}.

\begin{proof}[Filtrado]
Por la definición de probabilidad condicional y la de $\alpha_t$,
\[
\gammaf(k) = \Prob(z_t = k \mid x_{1:t})
= \frac{\Prob(x_{1:t}, z_t = k)}{\Prob(x_{1:t})}
= \frac{\alpha_t(k)}{\sum_{j=1}^{K} \alpha_t(j)},
\]
donde en el denominador se ha usado $\Prob(x_{1:t}) = \sum_j \Prob(x_{1:t}, z_t = j) = \sum_j \alpha_t(j)$.
\end{proof}

\begin{proof}[Causalidad]
Por la Proposición~\ref{prop:filtrado}, $\gammaf$ se expresa a través de $\{\alpha_t(j)\}_{j=1}^{K}$. La inicialización $\alpha_1(k) = \pi_k\, b_k(x_1)$ depende solo de $x_1$ y, por la recursión de la Proposición~\ref{prop:forward}, cada $\alpha_{t}(j)$ se obtiene de $\alpha_{t-1}(\cdot)$ y de $x_t$. Por inducción sobre $t$, $\alpha_t(j)$ es función de $x_{1:t}$, luego también lo es $\gammaf$.
\end{proof}

\section{Fórmulas de actualización del paso M de Baum-Welch}
\label{anx:em-pasos}

Desarrollo de las reestimaciones del paso M del algoritmo de Baum-Welch \cite{baum1970, rabiner1989}, Sección~\ref{sec:hmm-em}. La distribución inicial y la matriz de transición se actualizan como frecuencias esperadas,
\[
\hat{\pi}_k = \Prob(z_1 = k \mid x_{1:T}),
\qquad
\hat{A}_{ij} = \frac{\sum_{t=1}^{T-1} \Prob(z_t = i, z_{t+1} = j \mid x_{1:T})}{\sum_{t=1}^{T-1} \Prob(z_t = i \mid x_{1:T})},
\]
y las medias y covarianzas de cada régimen, como promedios ponderados por el posterior del estado,
\[
\hat{\mu}_k = \frac{\sum_{t=1}^{T} \Prob(z_t = k \mid x_{1:T})\, x_t}{\sum_{t=1}^{T} \Prob(z_t = k \mid x_{1:T})},
\qquad
\hat{\Sigma}_k = \frac{\sum_{t=1}^{T} \Prob(z_t = k \mid x_{1:T})\, (x_t - \hat{\mu}_k)(x_t - \hat{\mu}_k)^{\!\top}}{\sum_{t=1}^{T} \Prob(z_t = k \mid x_{1:T})},
\]
donde la media $\hat{\mu}_k$ que interviene en $\hat{\Sigma}_k$ es la ya reestimada en esa misma iteración.

\section{Recursión de Viterbi}
\label{anx:viterbi}

Desarrollo del algoritmo de Viterbi \cite{viterbi1967, rabiner1989}, Sección~\ref{sec:hmm-viterbi}. Si $\delta_t(j) = \max_{z_{1:t-1}} \Prob(z_{1:t-1},\, z_t = j,\, x_{1:t} \mid \lambda)$ es la probabilidad de la mejor trayectoria que termina en el estado $j$ en el instante $t$, la recursión es
\[
\delta_{t+1}(j) = b_j(x_{t+1}) \max_{1 \le k \le K} \delta_t(k)\, A_{kj},
\qquad \delta_1(k) = \pi_k\, b_k(x_1),
\]
y la trayectoria óptima se recupera memorizando en cada paso el estado $k$ que alcanza el máximo y recorriéndolos hacia atrás desde $T$.

\section{Monotonía del algoritmo EM}
\label{anx:em}

Demostración del Lema~\ref{lem:em} (Monotonía de EM), Sección~\ref{sec:hmm-em}, según el argumento clásico de EM \cite{dempster1977, bishop2006}.

\begin{proof}
Para cualquier parámetro $\theta$ descomponemos la log-verosimilitud observada tomando esperanza sobre $z_{1:T}$ con la ley del paso E, $\Prob(\cdot \mid x_{1:T}, \lambda)$, en la identidad $\log\Prob(x_{1:T}\mid\theta) = \log\Prob(x_{1:T}, z_{1:T}\mid\theta) - \log\Prob(z_{1:T}\mid x_{1:T}, \theta)$, cuyo miembro izquierdo no depende de $z_{1:T}$:
\[
\log\Prob(x_{1:T} \mid \theta) \;=\; Q(\theta \mid \lambda) - H(\theta \mid \lambda),
\]
con $Q(\theta\mid\lambda) = \E_{z\mid x,\lambda}\!\big[\log\Prob(x_{1:T}, z_{1:T}\mid\theta)\big]$ y $H(\theta\mid\lambda) = \E_{z\mid x,\lambda}\!\big[\log\Prob(z_{1:T}\mid x_{1:T},\theta)\big]$. El paso M elige $\lambda' = \arg\max_{\theta} Q(\theta\mid\lambda)$, luego $Q(\lambda'\mid\lambda) \ge Q(\lambda\mid\lambda)$. Por la desigualdad de Gibbs (un caso de la de Jensen),
\[
H(\lambda\mid\lambda) - H(\lambda'\mid\lambda) \;=\; \mathrm{KL}\big(\Prob(\cdot\mid x_{1:T},\lambda)\,\big\|\,\Prob(\cdot\mid x_{1:T},\lambda')\big) \;\ge\; 0.
\]
Restando ambas descomposiciones,
\[
\log\Prob(x_{1:T}\mid\lambda') - \log\Prob(x_{1:T}\mid\lambda)
= \underbrace{Q(\lambda'\mid\lambda) - Q(\lambda\mid\lambda)}_{\ge\, 0}
+ \underbrace{H(\lambda\mid\lambda) - H(\lambda'\mid\lambda)}_{\ge\, 0} \;\ge\; 0. \qedhere
\]
\end{proof}

\section{Estacionariedad y varianza incondicional del GARCH$(1,1)$}
\label{anx:garch-estac}

Demostración del Teorema~\ref{teo:garch-estac} (Estacionariedad en covarianza del GARCH$(1,1)$), Sección~\ref{sec:garch-estac}.

\begin{proof}
Supongamos primero que existe una solución estacionaria en covarianza, lo que significa, en particular, que $\Var(\epsilon_t)$ es finita y no depende de $t$; la llamamos $\bar{\sigma}^2 := \Var(\epsilon_t)$. Calculamos la varianza incondicional de $\epsilon_t$ por la ley de la varianza total, usando que $\E[\epsilon_t] = 0$ (pues $\E[\epsilon_t] = \E[\sigmat]\,\E[\eta_t] = 0$, al ser $\sigmat$ función del pasado e independiente de $\eta_t$, y $\E[\eta_t]=0$):
\[
\Var(\epsilon_t) = \E[\epsilon_t^2] = \E\!\big[\E[\epsilon_t^2 \mid \mathcal{F}_{t-1}]\big],
\]
donde $\mathcal{F}_{t-1}$ es la información hasta $t-1$. Como $\sigmat$ es función de $\mathcal{F}_{t-1}$ y $\eta_t$ es independiente de $\mathcal{F}_{t-1}$ con $\E[\eta_t^2] = \Var(\eta_t) = 1$,
\[
\E[\epsilon_t^2 \mid \mathcal{F}_{t-1}] = \sigmat^2\, \E[\eta_t^2 \mid \mathcal{F}_{t-1}] = \sigmat^2,
\]
de donde $\E[\epsilon_t^2] = \E[\sigmat^2]$. La misma identidad aplicada en $t-1$ da $\E[\epsilon_{t-1}^2] = \E[\sigma_{t-1}^2]$. Tomamos ahora esperanzas en la recursión \eqref{eq:garch11}:
\[
\E[\sigmat^2] = \omega + \alpha\, \E[\epsilon_{t-1}^2] + \beta\, \E[\sigma_{t-1}^2]
= \omega + \alpha\, \E[\sigma_{t-1}^2] + \beta\, \E[\sigma_{t-1}^2].
\]
Por estacionariedad, $\E[\sigmat^2] = \E[\sigma_{t-1}^2] =: m$, constante en el tiempo. Sustituyendo,
\[
m = \omega + (\alpha + \beta)\, m
\quad\Longleftrightarrow\quad
m\,(1 - \alpha - \beta) = \omega.
\]
Como $\omega > 0$ y $m = \E[\sigmat^2] = \Var(\epsilon_t) > 0$ es finita, el factor $1 - \alpha - \beta$ ha de ser estrictamente positivo, es decir, $\alpha + \beta < 1$, y entonces
\[
\Var(\epsilon_t) = m = \frac{\omega}{1 - \alpha - \beta}.
\]
Esto prueba la necesidad de la condición y la fórmula de la varianza.

Recíprocamente, supongamos $\alpha + \beta < 1$ y construyamos una solución estacionaria. Reescribiendo $\alpha \epsilon_{s-1}^2 + \beta \sigma_{s-1}^2 = (\alpha \eta_{s-1}^2 + \beta)\,\sigma_{s-1}^2$ e iterando la recursión \eqref{eq:garch11} hacia atrás $n$ pasos,
\[
\sigmat^2 = \omega \sum_{i=0}^{n-1} \prod_{l=1}^{i}\big(\alpha \eta_{t-l}^2 + \beta\big) \;+\; \Big(\prod_{l=1}^{n}\big(\alpha \eta_{t-l}^2 + \beta\big)\Big)\,\sigma_{t-n}^2.
\]
En lugar de suponer la solución dada, la \emph{construimos} como la serie
\[
\sigmat^2 \;:=\; \omega \sum_{i=0}^{\infty} \prod_{l=1}^{i}\big(\alpha \eta_{t-l}^2 + \beta\big),
\]
y comprobamos que está bien definida. Los $\eta_{t-l}$ son i.i.d. con $\E[\alpha \eta_{t-l}^2 + \beta] = \alpha + \beta$ y los factores de cada producto son independientes, de modo que el término $i$-ésimo tiene esperanza $\omega\,(\alpha+\beta)^i$. Como todos los sumandos son no negativos, el teorema de convergencia monótona permite intercambiar esperanza y suma:
\[
\E[\sigmat^2] \;=\; \omega \sum_{i=0}^{\infty} (\alpha+\beta)^i \;=\; \frac{\omega}{1-\alpha-\beta} \;<\; \infty,
\]
y la serie geométrica converge por ser $\alpha+\beta<1$. La serie que define $\sigmat^2$ converge así en $L^1$ (y casi seguramente), tiene esperanza finita y constante, y por construirse con las mismas variables i.i.d. desplazadas en el tiempo su ley no depende de $t$; además satisface la recursión \eqref{eq:garch11}. El término residual de la iteración finita, $\big(\prod_{l=1}^{n}(\alpha \eta_{t-l}^2+\beta)\big)\,\sigma_{t-n}^2$, tiene esperanza $(\alpha+\beta)^n\,\E[\sigma_{t-n}^2] \to 0$, lo que confirma que la serie es el límite de las iteraciones.

Queda ver que $\epsilon_t = \sigmat \eta_t$ es estacionario en covarianza, esto es, que su media, su varianza y sus autocovarianzas no dependen de $t$. La media es $\E[\epsilon_t] = \E[\sigmat]\,\E[\eta_t] = 0$, y la varianza, $\E[\epsilon_t^2] = \E[\sigmat^2] = \omega/(1-\alpha-\beta)$, es finita y constante. Para las autocovarianzas, fijado $h \ge 1$, como $\eta_t$ es independiente de $\mathcal{F}_{t-1}$ con $\E[\eta_t]=0$ y $\epsilon_{t-h}$ es $\mathcal{F}_{t-1}$-medible,
\[
\E[\epsilon_t\, \epsilon_{t-h}] = \E\!\big[\epsilon_{t-h}\,\E[\epsilon_t \mid \mathcal{F}_{t-1}]\big] = 0,
\qquad \text{ya que } \E[\epsilon_t \mid \mathcal{F}_{t-1}] = \sigmat\,\E[\eta_t] = 0.
\]
Media nula, varianza finita constante y autocovarianzas nulas: el proceso es estacionario en covarianza, lo que cierra la equivalencia \cite{nelson1990, francqzakoian2019}.
\end{proof}

\section{Verosimilitud condicional del GARCH$(1,1)$-$t$}
\label{anx:garch-mv}

Derivación de la log-verosimilitud condicional, Sección~\ref{sec:garch-mv}. Sea $f_\nu$ la densidad de la $t$ de Student estandarizada de $\nu$ grados de libertad. Condicionada a la información hasta $t-1$, la innovación $\epsilon_t = \sigmat \eta_t$ tiene densidad $\tfrac{1}{\sigmat} f_\nu(\epsilon_t/\sigmat)$, ya que es un cambio de escala de $\eta_t$ por el factor $\sigmat$, que es conocido en $t-1$. La log-verosimilitud condicional de la muestra $\epsilon_{1:T}$, dado un valor inicial $\sigma_1^2$, es la suma de las contribuciones diarias
\[
\ell(\theta) \;=\; \sum_{t=1}^{T} \left[\, \log f_\nu\!\left(\frac{\epsilon_t}{\sigmat}\right) - \log \sigmat \,\right],
\]
donde cada $\sigmat$ se obtiene recorriendo la recursión \eqref{eq:garch11} con los parámetros $\theta$ en evaluación. El término $-\log\sigmat$ es el jacobiano del cambio de variable de $\eta_t$ a $\epsilon_t$.

\section{Recursión de Adams--MacKay del BOCPD}
\label{anx:bocpd}

Demostración de la Proposición~\ref{prop:bocpd} (Recursión de Adams--MacKay) \cite{adams2007}, Sección~\ref{sec:bocpd-recursion}.

\begin{proof}
Partimos de la conjunta del instante $t$ e introducimos la run-length anterior marginalizando sobre sus valores posibles:
\[
\Prob(r_t, x_{1:t}) = \sum_{r_{t-1}} \Prob(r_t, r_{t-1}, x_{1:t}).
\]
Cada sumando se descompone separando la nueva observación $x_t$ del pasado $x_{1:t-1}$ mediante la regla del producto,
\[
\Prob(r_t, r_{t-1}, x_{1:t})
= \Prob(r_t \mid r_{t-1}, x_{1:t-1})\,
  \Prob(x_t \mid r_t, r_{t-1}, x_{1:t-1})\,
  \Prob(r_{t-1}, x_{1:t-1}).
\]
El primer factor es la dinámica de la run-length, que por construcción del modelo solo depende de $r_{t-1}$: dado un tramo de longitud $r_{t-1}$, en el paso siguiente o bien hay cambio, y entonces $r_t = 0$ con probabilidad $H(r_{t-1}+1)$, o bien no lo hay, y entonces $r_t = r_{t-1}+1$ con probabilidad $1 - H(r_{t-1}+1)$. Es decir,
\[
\Prob(r_t \mid r_{t-1}, x_{1:t-1}) =
\begin{cases}
H(r_{t-1}+1), & r_t = 0,\\[2pt]
1 - H(r_{t-1}+1), & r_t = r_{t-1}+1,\\[2pt]
0, & \text{en otro caso.}
\end{cases}
\]
El segundo factor es la predictiva del dato nuevo dentro del tramo: condicionada a $r_{t-1}$, y por la independencia entre segmentos, la observación $x_t$ depende solo de las $r_{t-1}$ observaciones del tramo en curso, lo que da $\pi^{(r_{t-1})}_{\mathrm{pred}}$. El tercer factor es la conjunta del paso anterior, ya disponible. En ambos mensajes interviene esta misma predictiva $\pi^{(r_{t-1})}_{\mathrm{pred}}$ porque $x_t$ se evalúa contra el segmento en curso, de longitud $r_{t-1}$: el reinicio se materializa en la run-length siguiente, $r_t = 0$, no en la pertenencia de $x_t$, que todavía corresponde a ese segmento.

Como la transición de la run-length concentra toda su masa en dos destinos, la marginalización sobre $r_{t-1}$ se reparte en dos casos. Si $r_t = r_{t-1} + 1 > 0$, solo contribuye el valor concreto $r_{t-1} = r_t - 1$ (no hay otra forma de llegar a esa run-length sin reinicio), de donde
\[
\Prob(r_t = r_{t-1}+1,\, x_{1:t})
= \Prob(r_{t-1}, x_{1:t-1})\,\big(1 - H(r_{t-1}+1)\big)\,\pi^{(r_{t-1})}_{\mathrm{pred}},
\]
que es el mensaje de crecimiento. Si $r_t = 0$, contribuyen \emph{todos} los valores de $r_{t-1}$, pues desde cualquier longitud puede producirse un reinicio, y la suma da
\[
\Prob(r_t = 0,\, x_{1:t})
= \sum_{r_{t-1}} \Prob(r_{t-1}, x_{1:t-1})\, H(r_{t-1}+1)\, \pi^{(r_{t-1})}_{\mathrm{pred}},
\]
que es el mensaje de reinicio. Finalmente, la posterior se obtiene de la conjunta dividiendo por la verosimilitud marginal del dato, $\Prob(x_{1:t}) = \sum_{r}\Prob(r_t = r, x_{1:t})$, lo que da la normalización del enunciado.
\end{proof}

\section{Pseudo-residuos del boosting logístico}
\label{anx:gbm-residuo}

Enunciado y demostración del resultado citado en la Sección~\ref{sec:sup-gbm} (gradient boosting).

\begin{proposicion}[Pseudo-residuos del boosting logístico]
\label{prop:gbm-residuo}
Si $F$ es el logaritmo de la razón de momios y $p = \sigma(F) = (1 + e^{-F})^{-1}$ la probabilidad asociada, el pseudo-residuo de la pérdida logística es la diferencia entre la etiqueta y la probabilidad predicha,
\[
r_i = -\frac{\partial L}{\partial F}\bigg|_{F(\mathbf{x}_i)} = y_i - p(\mathbf{x}_i).
\]
\end{proposicion}

\begin{proof}
Con $p = \sigma(F)$ se tiene $\mathrm{d}p/\mathrm{d}F = p(1-p)$. La contribución de un dato a la pérdida es $L = -y\log p - (1-y)\log(1-p)$, de donde
\[
\frac{\partial L}{\partial p} = -\frac{y}{p} + \frac{1-y}{1-p} = \frac{p - y}{p(1-p)}.
\]
Por la regla de la cadena, $\partial L/\partial F = \big(\partial L/\partial p\big)\,p(1-p) = p - y$, y el pseudo-residuo es su opuesto, $y - p$ \cite{friedman2001}.
\end{proof}

\section{Pesos óptimos y ganancia de una división en XGBoost}
\label{anx:xgb}

Enunciado y demostración del resultado citado en la Sección~\ref{sec:sup-xgboost} (XGBoost).

\begin{proposicion}[Pesos óptimos y ganancia de una división]
\label{prop:xgb-score}
Fijada la estructura $q$, con conjuntos de instancias por hoja $I_j = \{i : q(\mathbf{x}_i) = j\}$ y sumas $G_j = \sum_{i \in I_j} g_i$, $H_j = \sum_{i \in I_j} h_i$, el peso que minimiza la aproximación de segundo orden del objetivo y el valor óptimo correspondiente son
\[
w_j^{\star} = -\frac{G_j}{H_j + \lambda},
\qquad
\widetilde{\mathcal{L}}^{\star}(q) = -\frac{1}{2}\sum_{j=1}^{T} \frac{G_j^2}{H_j + \lambda} + \gamma\, T.
\]
En consecuencia, dividir una hoja en dos descendientes $L$ y $R$ reduce el objetivo en
\[
\mathrm{Gan} = \frac{1}{2}\left[\frac{G_L^2}{H_L + \lambda} + \frac{G_R^2}{H_R + \lambda} - \frac{(G_L + G_R)^2}{H_L + H_R + \lambda}\right] - \gamma,
\]
que es el criterio con el que se eligen los cortes.
\end{proposicion}

\begin{proof}
La expansión de Taylor de segundo orden de la pérdida y el descarte del término constante $l(y_i, \widehat{y}_i^{(t-1)})$ dan
\[
\widetilde{\mathcal{L}}^{(t)} = \sum_{i=1}^{n} \big[\, g_i f_t(\mathbf{x}_i) + \tfrac{1}{2} h_i f_t(\mathbf{x}_i)^2 \,\big] + \Omega(f_t).
\]
Como $f_t(\mathbf{x}_i) = w_j$ para todo $i \in I_j$, se agrupa por hoja,
\[
\widetilde{\mathcal{L}}^{(t)} = \sum_{j=1}^{T} \Big[\, G_j\, w_j + \tfrac{1}{2}(H_j + \lambda)\, w_j^2 \,\Big] + \gamma\, T.
\]
Para la pérdida logística, la Proposición~\ref{prop:gbm-residuo} da $g_i = p_i - y_i$ y un cálculo análogo $h_i = p_i(1 - p_i) \ge 0$, luego $H_j + \lambda > 0$ y cada sumando es una parábola convexa en $w_j$. Su mínimo está en $w_j^{\star} = -G_j/(H_j + \lambda)$, con valor $-\tfrac{1}{2} G_j^2/(H_j + \lambda)$; sumando sobre las hojas se obtiene $\widetilde{\mathcal{L}}^{\star}(q)$. La ganancia de una división es la diferencia entre el objetivo de la hoja madre y el de los dos hijos, que arroja la expresión del enunciado \cite{chen2016}.
\end{proof}
